\chapter{DP-600 Exam Mapping and Study Guide}
\index{DP-600!exam mapping}\index{certification!study guide}%
This appendix maps the DP-600 \emph{Implementing Analytics Solutions Using Microsoft Fabric}
exam to the chapters of this book, then turns that mapping into a concrete study strategy and
self-assessment checklist. Chapter~24 runs this process as a guided week; this appendix is the
standalone reference version. Always verify domains and weights against the live skills
outline---it evolves as Fabric evolves \cite{microsoftDp600}.

\section{Exam Profile}
\index{DP-600!profile}%
\begin{itemize}
    \item \textbf{Audience.} Analytics engineers who design, build, and operate Fabric
    analytics solutions: lakehouses, warehouses, pipelines, real-time intelligence, and
    semantic models.
    \item \textbf{Format.} Scenario-based multiple choice plus case studies; expect
    ``which technology / which setting / what happens next'' questions, not definitions.
    \item \textbf{Hands-on value.} High. The exam's scenarios mirror this book's labs---Direct
    Lake fallback triage, gateway placement, RLS verification, capacity throttling---which is
    why remediation below means \emph{rebuilding}, not rereading.
    \item \textbf{Adjacent certifications.} PL-300 (Power BI Data Analyst) for report-first
    careers; the DP-203 lineage for Azure-data-engineering context.
\end{itemize}

\section{Domain-to-Chapter Map}
\index{DP-600!domain map}%
The skills outline groups objectives into weighted domains. Table~\ref{tab:appendix-e-domains}
maps each to the chapters that teach it; the paraphrase follows Chapter~24.

\begin{table}[H]
\centering
\small
\begin{tabular}{@{}p{5.8cm}p{4.2cm}p{3.6cm}@{}}
\toprule
\textbf{Skills area (paraphrased)} & \textbf{This book} & \textbf{Typical weight} \\
\midrule
Plan, implement, and manage an analytics solution (capacity, workspace, governance, CI/CD) & Chapters 1, 21--23 & Largest single share \\
Prepare and serve data (lakehouse, warehouse, pipelines, dataflows, notebooks) & Chapters 2--5, 9--12 & Large \\
Implement and manage semantic models (modeling, DAX, Direct Lake) & Chapters 6--8 & Large \\
Real-time intelligence (eventstreams, KQL, activator) & Chapters 13--16 & Moderate and growing \\
Data science and ML touchpoints & Chapters 17--20 & Smaller but present \\
\bottomrule
\end{tabular}
\caption{DP-600 domains mapped to chapters (from Chapter~24). Verify current weights against
the live skills outline \cite{microsoftDp600}.}
\label{tab:appendix-e-domains}
\end{table}

\section{Objective-Level Mapping}
\index{DP-600!objectives map}%
Within each domain, the objectives below are the ones candidates most often under-prepare.
Each entry names the chapter section whose lab builds the tested judgment.

\subsection{Plan, Implement, and Manage an Analytics Solution}
\begin{itemize}
    \item Capacity SKUs, CU metering, smoothing and throttling behavior, surge protection ---
    Ch.~1 (capacities), Ch.~23 (Capacity Metrics app, throttling stages).
    \item Workspace design, domains, roles, and item-level sharing --- Ch.~1, Ch.~21.
    \item Git integration, deployment pipelines, and deployment rules --- Ch.~23.
    \item Sensitivity labels, endorsement, lineage, and Purview integration --- Ch.~22.
    \item Disaster recovery expectations: Git rehydration, Delta time travel, no automatic
    cross-region failover --- Ch.~23.
\end{itemize}

\subsection{Prepare and Serve Data}
\begin{itemize}
    \item Lakehouse versus Warehouse decision criteria --- Ch.~2, Ch.~5.
    \item Shortcuts versus copy/materialization --- Ch.~1.
    \item Spark pools, environments, notebooks, and Delta optimization
    (\texttt{OPTIMIZE}, V-Order, Z-Order, \texttt{VACUUM}, time travel) --- Ch.~3, 4.
    \item \texttt{COPY INTO}, T-SQL transformation, three-part cross-item queries --- Ch.~5.
    \item Pipeline activities (Lookup, ForEach, If, Until), parameterization, scheduling ---
    Ch.~9.
    \item Dataflows Gen2 destinations, staging, incremental refresh --- Ch.~10.
    \item On-premises and VNet data gateways --- Ch.~11.
    \item Orchestration patterns, Monitoring Hub, alerting --- Ch.~12.
\end{itemize}

\subsection{Implement and Manage Semantic Models}
\begin{itemize}
    \item Star-schema modeling, relationships, bridge tables, SCD handling --- Ch.~7.
    \item DAX measures and calculation discipline --- Ch.~7.
    \item Storage modes: Import, DirectQuery, Direct Lake; framing; fallback triggers and their
    telemetry --- Ch.~8.
    \item SQL analytics endpoint objects: views, TVFs, RLS, column-level security --- Ch.~6.
\end{itemize}

\subsection{Real-Time Intelligence}
\begin{itemize}
    \item Eventstream sources, destinations, and fan-out topology --- Ch.~13.
    \item Eventhouses, KQL databases, materialized views, dedup idioms --- Ch.~14.
    \item Data Activator (Reflex) triggers from streams and Power BI visuals --- Ch.~15.
    \item Real-Time Dashboards and automatic page refresh --- Ch.~16.
\end{itemize}

\subsection{Data Science and ML Touchpoints}
\begin{itemize}
    \item Notebooks, Data Wrangler, and feature tables --- Ch.~17.
    \item Built-in MLflow experiments and the model registry --- Ch.~18.
    \item SynapseML and Azure AI service integration --- Ch.~19.
    \item Batch scoring with \texttt{PREDICT} and pipeline-orchestrated inference --- Ch.~20.
\end{itemize}

\section{Study Strategy}
\index{DP-600!study strategy}%
\begin{enumerate}
    \item \textbf{Build the readiness matrix first.} One row per domain, columns for evidence
    (which chapter labs you completed, which capstone artifacts exist) and confidence. Domains
    with neither are the plan. (Ch.~24, Monday.)
    \item \textbf{Study to the outline, not the table of contents.} Re-download the skills
    outline in exam week and diff it against your matrix; Fabric evolves faster than any book.
    \item \textbf{Take two timed, closed-book mock exams} scored by domain. Classify every miss
    as knowledge (rebuild the lab), misread (fix pacing), or judgment (re-derive the Friday use
    cases). (Ch.~24, Wednesday/Friday.)
    \item \textbf{Remediate weakest domain first, by building.} Recognition memory from
    rereading produces confidence without capability; the exam tests scenario decisions.
    \item \textbf{Memorize the decision tables}: Lakehouse vs.\ Warehouse (Ch.~2, 5), Import vs.\
    Direct Lake vs.\ DirectQuery (Ch.~8), gateway types (Ch.~11), Eventstream destinations
    (Ch.~13), batch vs.\ streaming (Ch.~13, 20).
    \item \textbf{Use the readiness threshold}: two consecutive full-length mocks at 80\%+ with
    no domain below 70\%, plus the ability to explain why each wrong option is wrong.
\end{enumerate}

\section{Self-Assessment Checklist}
\index{DP-600!self-assessment}%
You are ready when you can honestly tick every line---with an artifact or a rebuilt lab as
evidence for each.

\subsection*{Platform and Administration}
\begin{itemize}
    \item[$\square$] I can explain CU metering, smoothing, and the throttling stages, and read
    them in the Capacity Metrics app. (Ch.~1, 23)
    \item[$\square$] I can design a dev/test/prod workspace topology with Git binding and
    deployment rules. (Ch.~23)
    \item[$\square$] I can state what Git rehydration does and does not recover. (Ch.~23)
\end{itemize}

\subsection*{Data Preparation and Serving}
\begin{itemize}
    \item[$\square$] I can choose Lakehouse vs.\ Warehouse for a stated workload and defend the
    choice. (Ch.~2, 5)
    \item[$\square$] I can write and troubleshoot \texttt{COPY INTO}, including OneLake path
    forms. (Ch.~5, Appendix~B)
    \item[$\square$] I can schedule \texttt{OPTIMIZE}/V-Order and explain \texttt{VACUUM}
    retention risk. (Ch.~4)
    \item[$\square$] I can build a metadata-driven pipeline with a contract gate. (Ch.~9)
    \item[$\square$] I can pick the correct gateway type for an on-prem or VNet source. (Ch.~11)
    \item[$\square$] I can wire On Failure alerting with a runbook link. (Ch.~12)
\end{itemize}

\subsection*{Semantic Models}
\begin{itemize}
    \item[$\square$] I can model a star schema with declared grain and conformed dimensions. (Ch.~7)
    \item[$\square$] I can name Direct Lake fallback triggers and find them in telemetry. (Ch.~8)
    \item[$\square$] I can implement and test RLS as each persona. (Ch.~6, 21)
\end{itemize}

\subsection*{Real-Time Intelligence}
\begin{itemize}
    \item[$\square$] I can route an eventstream to KQL and lakehouse destinations and justify
    the fan-out. (Ch.~13)
    \item[$\square$] I can write the \texttt{arg\_max} dedup idiom and a materialized view. (Ch.~14)
    \item[$\square$] I can configure a sustained-threshold Reflex alert that fires once. (Ch.~15)
\end{itemize}

\subsection*{ML Touchpoints}
\begin{itemize}
    \item[$\square$] I can track an experiment and load a model by registry alias. (Ch.~18)
    \item[$\square$] I can explain where SynapseML fits and how credentials are vaulted. (Ch.~19)
    \item[$\square$] I can describe batch scoring with \texttt{model\_version} auditing. (Ch.~20)
\end{itemize}

\begin{pitfall}
The most common DP-600 failure pattern is over-weighting the domains you already know. The
readiness matrix exists to prevent exactly that: let evidence and confidence---not
comfort---set the study order.
\end{pitfall}
